% !TeX root = ../main.tex
\newpage
\section{Conclusiones}
En esta práctica se ha obtenido una ley empírica que relaciona la longitud crítica de un transistor con la masa efectiva del material en su canal, bajo condiciones específicas de una barrera de potencial de $0.5$\,eV y una temperatura de $300$\,K. Se ha establecido que esta ley, descrita por la ecuación \eqref{eq:lym*}, puede generalizarse a otras condiciones siempre que se mantenga la dependencia fundamental entre la longitud crítica y la masa efectiva.

Los resultados sugieren que el modelo propuesto se adecua a la realidad física, ofreciendo una visión clave sobre los límites de la miniaturización de transistores tipo MOSFET. Aunque este estudio no representa una solución exhaustiva al problema de la miniaturización, aborda un aspecto fundamental: la relación intrínseca entre la masa efectiva del portador y la dimensión mínima del canal. La comprensión de esta relación es crucial para el diseño futuro de transistores más pequeños y eficientes.

Las limitaciones de este estudio incluyen la simplificación del modelo a un sistema 1D, lo que puede no capturar completamente los efectos tridimensionales presentes en dispositivos reales. Además, la ley empírica obtenida se basa en un rango específico de masas efectivas y condiciones operativas, por lo que su aplicabilidad a otros materiales o configuraciones podría requerir ajustes adicionales. Sin embargo, la alta correlación obtenida en el ajuste de la ley de potencias sugiere que la relación propuesta es robusta dentro del rango estudiado.\\

Otro elemento importante que no se ha tenido en cuenta en este estudio es la eficiencia de los transistores de alta masa efectiva. Aunque preveemos que los transistores de masa efectiva mayor tienen un tamaño menor, no se ha analizado si esta reducción de tamaño se traduce en una mejora de eficiencia para el dispositivo. Además, cabe esperar que al ofrecer una movilidad menor, los transistores de alta masa efectiva den lugar a una mayor disipación de energía térmica, lo que podría limitar su aplicabilidad en ciertos contextos. Por tanto, sería interesante analizar cómo varía la eficiencia de los transistores a medida que se reduce la longitud del canal, y cómo esta eficiencia se relaciona con la masa efectiva del material utilizado.\\

En resumen, este trabajo ha permitido validar experimentalmente la dependencia entre la longitud crítica y la masa efectiva. Como posible línea de investigación futura, sería interesante analizar cómo la variación de la altura de la barrera de potencial o de la temperatura afecta a la ley empírica obtenida, permitiendo así refinar el modelo para un rango más amplio de condiciones operativas.\\

Y por supuesto, un último comentario sobre la ley de Moore y por qué da asco: la ley de Moore fue propuesta por el cofundador de Intel, Gordon Moore, en 1965. Además de empresario el chaval pues hace pila de tiempo de eso, abe?. Y aunque ha funcionado durante décadas, la ley de Moore es una idealización del avance tecnológico y no es una ley física fundamental. Así que sí, es una mierda, pero ha guiado a muchas industrias durante muchos años. En definitiva, con ley de Moore o sin ella, el avance de la ciencia y de la tecnología seguirá su curso, aunque quizás a un ritmo diferente al que hemos visto en el pasado.